\newpage
\section{Revisão bibliográfica}
\label{sc:revisao}





A seção de revisão bibliográfica tem como objetivo apresentar e discutir os principais estudos e contribuições já realizados sobre o tema, organizando o conhecimento existente de forma lógica e crítica. Para preenchê-la:

\begin{enumerate}
    \item \textbf{Organize o conteúdo}: Estruture a revisão por temas, abordagens, cronologia ou outras categorias que facilitem a compreensão do leitor.
    \item \textbf{Sintetize os principais trabalhos}: Destaque os estudos mais relevantes, mencionando seus autores, métodos, principais achados e limitações.
    \item \textbf{Identifique lacunas}: Aponte aspectos ainda pouco explorados ou controversos, enfatizando oportunidades para contribuições futuras.
    \item \textbf{Relacione com seu tema}: Explique como os estudos revisados estão conectados ao objetivo do seu trabalho e de que forma embasam sua discussão.
    \item \textbf{Use referências de qualidade}: Priorize fontes confiáveis, como artigos científicos, livros e relatórios reconhecidos na área.
\end{enumerate}

Mantenha o texto coeso e evite apenas listar trabalhos; ao invés disso, discuta as contribuições e conexões entre eles para oferecer uma visão abrangente e crítica do estado da arte.
